\chapter{Podsumowanie}

Celem pracy była implementacja oraz analiza wydajności algorytmu aktywnego
uczenia się deterministycznych automatów visibly pushdown. W ramach pracy
przedstawiono podstawy teoretyczne dotyczące automatów DVPA oraz aktywnego
uczenia automatów, a następnie opisano zaimplementowany algorytm ucznia i
nauczyciela. Omówiono również najważniejsze decyzje implementacyjne oraz moduły
projektu umożliwiające testowanie algorytmu na przygotowanych klasach automatów.

Drugim istotnym elementem pracy było rozszerzenie algorytmu o wykorzystanie
więzów opartych na systemach przepisywania słów. W pracy opisano sposób
adaptacji tej idei do rozważanego modelu automatów DVPA oraz przedstawiono
implementację mechanizmu sprawdzania zgodności hipotezy z zadanym systemem
SRS. Przygotowane eksperymenty pozwoliły zbadać wpływ tej funkcjonalności na
liczbę zadawanych zapytań o równoważność oraz na skuteczność uczenia w przypadku
heurystycznego sprawdzania równoważności.

Przeprowadzone eksperymenty pokazały, że największy wpływ na liczbę zadawanych
zapytań ma liczba stanów automatu. Zmiany rozmiaru alfabetu wejściowego oraz
alfabetu stosowego miały mniejsze znaczenie, przy czym spośród analizowanych
typów symboli największy wpływ na wydajność algorytmu miała liczba symboli typu
return. Zaobserwowano również, że średnia liczba zapytań o równoważność była w~przybliżeniu
równa liczbie stanów uczonego automatu.

Eksperymenty z systemem SRS pokazały, że wykorzystanie dodatkowych podpowiedzi
ogranicza liczbę kosztownych zapytań o równoważność. Może również prowadzić do
skrócenia czasu uczenia, szczególnie dla odpowiednio dużych automatów. Dodatkowo może
wpływać na skuteczność uczenia przy zastosowaniu heurystycznych metod
sprawdzania równoważności, przy czym uzyskana poprawa w dużej mierze zależy
od konstrukcji zastosowanych reguł.

Należy zaznaczyć, że przeprowadzona analiza miała charakter eksperymentalny i
była ograniczona do przygotowanych klas generatorów oraz wybranych konfiguracji
parametrów. Możliwym kierunkiem dalszego rozwoju jest rozszerzenie zbioru
generatorów, w szczególności w kontekście badania różnych typów reguł SRS oraz
ich wpływu na wydajność uczenia. Interesującym zagadnieniem implementacyjnym
pozostaje również poszukiwanie dalszych optymalizacji w obrębie algorytmu
nauczyciela. Powstała architektura może także służyć jako punkt wyjścia do
porównywania oraz testowania heurystycznych algorytmów sprawdzania
równoważności w algorytmach aktywnego uczenia.